\newpage
\begin{appendix}

\renewcommand{\thechapter}{\Alph{chapter}}
\setcounter{figure}{0}
\setcounter{table}{0}
\setcounter{equation}{0}
\counterwithin{figure}{chapter}
\counterwithin{table}{chapter}
\counterwithin{equation}{chapter}
\renewcommand{\thefigure}{\thechapter-\arabic{figure}}
\renewcommand{\thetable}{\thechapter-\arabic{table}}
\renewcommand{\theequation}{\thechapter-\arabic{equation}}

\chapter{}
附录的有无，根据毕业设计（论文）情况而定，内容一般包括正文内不便列出的冗长公式推导、辅助性数学工具、符号说明（含缩写）、计算程序及说明等。附录另起一页。附录依序用大写正体A、B、C…编序号，如：附录A。附录中的图、表、公式、参考文献等另行编序号，与正文分开，也一律用阿拉伯数字编码，但在数码前冠以附录序码，如：图A01、表B-2、公式（B-3）、文献[A5]等。“附录”两字之间空一个全角空格或两个半角空格，采用不编号章标题样式；内容应采用正文样式。

这里援引一个\cref{eq:emc2_1}

\begin{equation}
  E = mc^2
  \label{eq:emc2_1}
\end{equation}

\end{appendix}